\documentclass[11pt]{report}
\usepackage[margin=0.59in]{geometry}
\usepackage{graphicx}
\usepackage{tikz}
\usetikzlibrary{calc}
\usepackage{hyperref}


\newcommand{\mediumsize}{\fontsize{15pt}{16pt}\selectfont} 
\begin{document} 
\begin{titlepage}

    \begin{minipage}[t]{0.95\textwidth}
        \hfill
        \raggedleft
        \textbf{Immanuel Dietrich} \\
        Utrecht, Netherlands \\
        \today\\
        \href{mailto:immanueld49@gmail.com}{immanueld49@gmail.com} \\
        +49 1512 9081675 \\
        \url{http://linkedin.com/in/immanuel-dietrich}
    \end{minipage}

\raggedright \textbf{Rijkswaterstaat}  \\  Griffioenlaan 2 \\  3526 LA Utrecht \\ 

\vspace{0.7em}

\raggedright Beste lezer,\\

\vspace{0.7em}

Met deze motivatiebrief wil ik graag solliciteren naar de positie van Adviseur chemische waterkwaliteit bij Rijkswaterstaat. De werkzaamheden en doelstellingen van deze functie spreken mij bijzonder aan. 
Hoewel ik inmiddels bijna twee jaar in Nederland (Utrecht) woon, ben ik de taal nog aan het leren en druk ik mijzelf makkelijker uit in het Engels:

\vspace{0.7em}

\section*{1. Motivation}

Growing up partially in the Salzkammergut,  a lake district in the Austrian Alps, I was captivated by how dramatically the water cycle shifted between seasons - Ice melting into roaring rivers, lakes transforming with snowmelt. That early fascination never left me. 
During my Erasmus semester in Vienna, I met a PhD researcher working on reverse osmosis technology. Through her, I first encountered the Netherlands' approach to water management - not just as infrastructure, but as a national identity built around living with water. 
That conversation was a turning point. I chose to pursue my Master's in Water Science and Management in Utrecht specifically because of it, and after two years here I am certain this is where I want to build my career.

What draws me to this role at Rijkswaterstaat specifically is the combination of chemical water quality assessment, policy translation, and multi-stakeholder work. My academic background spans PFAS, EU water legislation, 
the Water Framework Directive, and drinking water treatment, but I am equally motivated by the applied side. Turning scientific analysis into practical advice that protects our water systems. Rijkswaterstaat operates at exactly that space between science, 
regulation, and impact, which is where I want to work.

\vspace{1em}

\section*{2. Organisational Sensitivity}

\textbf{Situation:} During my bachelor's thesis internship, I worked on an EU-funded drinking water treatment project in Kikinda, Serbia. As one of the few international team members, I collaborated with local utility employees operating under different 
organisational norms and priorities.

\textbf{Task:} Besides carrying out my technical work, I needed to coordinate with local colleagues to ensure that the project progressed despite differences in working methods, priorities and communication styles. 
My technical work involved evaluating coagulation and flocculation parameters for a new treatment plant.

\textbf{Action:} Initially I pushed my own approach, which only increased resistance. I stepped back, listened to their concerns, and reframed discussions about our shared objective: delivering clean drinking water to the city. I adapted how I communicated. 
I chose for less directive, more collaborative and made sure colleagues understood why specific steps mattered rather than just what needed to be done. 

\textbf{Result:} The project was a success. The plant was successfully commissioned and brought drinking water to a community that had previously been supplied with water violating Serbian legal standards.

\textbf{Reflection:} I learned that in politically or culturally sensitive environments, listening and adapting your communication style are prerequisites for technical progress, not optional extras.

\vspace{1em}

\section*{3. Content Knowledge}


\textbf{Situation:} DDuring my Master's, I contributed to a consultancy project for KWR on adaptive strategies for the Dutch drinking water system under deep uncertainty. My work focused on assessing vulnerabilities in current treatment and distribution 
infrastructure and developing solution pathways for future scenarios.

\textbf{Task:} My objective was to evaluate the technological readiness of molecularly imprinted polymers for PFAS removal and assess whether they could become a viable treatment technology.

\textbf{Action:} As a team, we analysed the Dutch drinking water system across its full chain identifying structural weaknesses such as dependency on specific source waters, limited treatment flexibility, and regulatory gaps in emerging contaminant management. 
Drawing on the Water Framework Directive and Dutch drinking water legislation, we evaluated where the system was insufficiently prepared for pressures including climate change, new pollutants, and demand shifts.

\textbf{Result:} Based on this analysis, we developed concrete recommendations covering treatment upgrades, source diversification, and adaptive governance measures. We were translating technical and regulatory knowledge into practical, prioritised options for 
decision-makers.

\textbf{Reflection:} This project required me to integrate scientific literature, legislative frameworks, and practical system knowledge simultaneously. It strengthened my ability to produce structured, evidence-based advice under uncertainty, 
which I see as central to the advisory work at Rijkswaterstaat, where chemical water quality decisions must be grounded in both technical rigour and regulatory context.

\vspace{1em}

\section*{4. Analysing}

\textbf{Situation:} During my Master's thesis on PFAS removal technologies, I had to process a large body of scientific literature alongside my own experimental data. It covered competing technologies, varying performance metrics, and often contradictory findings.

\textbf{Task:} My approach was to first build structure before attempting analysis. I categorised studies by technology type, then broke each down by key parameters: removal efficiency, selectivity, regeneration potential and scalability. 
I built comparison tables and flow diagrams to make relationships and gaps visible, which prevented me from getting lost in volume and helped me spot inconsistencies across studies


\textbf{Action:} Throughout the process I kept returning to my core research question - Is this piece of information actually relevant to my goal? That helped me to avoid getting sidetracked by further information that might be interesting, 
but less relevant for my research. 

\textbf{Result:} This approach produced a clear technological readiness assessment with a structured barrier analysis, giving my supervisors and readers a usable framework rather than an information dump.

\textbf{Reflection:} I apply the same logic outside research: organise first, analyse second, and continuously test whether what you're looking at serves the actual question. In an advisory context, that discipline is what separates useful analysis from noise.

Thank you for considering my application.

\vspace{0.7em}

\raggedright Kind regards,\\
\textbf{Immanuel Dietrich}

\end{titlepage}
\end{document}
